\subsection{The simple tangent-triple type for every translated triple}
\label{sec:DE-A4-uniform-tangent-triple-type}

\begin{lemma}[A uniform local type on the main triple-contact incidence]
\label{lem:DE-A4-uniform-triple-local-type}
\status{Proved. Complete independent proof and program audits by root
and monodromy passed.}
Let $E:y^2+y=x^3$ over an algebraically closed field of
characteristic two, $X=E^2$, and
$L=\OO_E(3O)\boxtimes\OO_E(3O)$.
Let $K\subset X(k)$ be a finite odd subgroup with injective
coordinate projections.  Fix $0\ne g\in K$ and
$\eta\in K\setminus\{0,g\}$.
Consider sections containing $R,R+g,R+\eta$, with a common translated
tangent at the first two points.

The locus where these five linear conditions are independent is
an irreducible nonempty open in the base of marked points and
tangent directions.  Over that open the coefficient incidence is
a $\mathbf P^3$-bundle.  It contains a nonempty open of sections
whose source curve is globally smooth, satisfies
$\det S\ne0$ and $Q_H(S)\ne0$, has simple contact of order exactly
two at the first two translated branches, and has a transverse
third branch.
There is no exception when $g$ has order three or when
$\{0,g,\eta\}$ is a cyclic three-element subgroup.
\end{lemma}
\begin{proof}
Write $h_0=0,h_1=g,h_2=\eta$, and let
\[
                    p_i=R_1+(h_i)_1,\qquad
                    q_i=R_2+(h_i)_2
                    \quad(0\le i\le2).
\]
All three $p_i$ are distinct, as are all three $q_i$.
Choose $R$ in a dense open where each triple of value vectors
in the plane-cubic embedding is a basis.
Such a choice is possible: three points on this cubic are
collinear exactly when their sum is zero, so the failures are
the finite conditions $3R_j+g_j+\eta_j=0$ on each elliptic factor.
Normalize the value vectors to $e_0,e_1,e_2$ separately in the two
factors.  A section is then represented by a matrix $S$.
Its containment of the three points says
\[
                             S_{00}=S_{11}=S_{22}=0.
\]

Use the invariant formal-group parameters on both factors.
They identify the tangent directions at all translated points.
Write the first-factor tangent vectors in the normalized frames as
$a_i=(a_{i0},a_{i1},a_{i2})$, and the second-factor tangent vectors
as $b_i=(b_{i0},b_{i1},b_{i2})$.
We may also require
\[
             a_{02}a_{12}b_{02}b_{12}b_{20}b_{21}\ne0,
             \qquad q_0,q_1\notin E[3].
\]
Every excluded condition is finite in the corresponding $R_j$.
For example, $a_{02}=0$ says that the tangent at $p_0$ contains
$p_1$, equivalently $3R_1+g_1=0$.
The condition $a_{12}=0$ is $3R_1+2g_1=0$.
The analogous conditions for $b_{02},b_{12}$ use the second
factor, and those for $b_{20},b_{21}$ are
$3R_2+2\eta_2+g_2=0$ and $3R_2+2\eta_2=0$.
The nonflex conditions are again finite.  Multiplication by three
is etale and finite in characteristic two, so these exclusions
cannot exhaust either elliptic curve.

Take the common tangent direction to be the second-factor direction.
After the three value conditions, the two tangent conditions
are independent equations on rows zero and one.
All sections satisfying these five conditions therefore have the
form
\[
 S=
 \begin{pmatrix}
       0&\alpha b_{02}&\alpha b_{01}\\
       \beta b_{12}&0&\beta b_{10}\\
       r&s&0
 \end{pmatrix}.
\]
Here $\alpha,\beta,r,s$ are free.  Fix $\alpha\beta\ne0$.
The row-zero and row-one sections on the second cubic are
nonzero multiples of its tangent-line sections at $q_0,q_1$.
Because these points are not flexes, their local expansions have
nonzero quadratic coefficients $\alpha\kappa_0$ and
$\beta\kappa_1$, with $\kappa_0\kappa_1\ne0$.

The first-factor derivatives of the full source equation at the
first two points are
\[
        l_0=a_{01}\beta b_{12}+a_{02}r,\qquad
        l_1=a_{10}\alpha b_{02}+a_{12}s.
\]
They are independent affine functions of $r,s$.
Choose them nonzero, so the source is smooth at both points and
has the indicated common tangent.
In the common invariant formal parameters $(u,t)$, the two
branches then have expansions
\[
       u=\frac{\alpha\kappa_0}{l_0}t^2+O(t^3),
       \qquad
       u=\frac{\beta\kappa_1}{l_1}t^2+O(t^3).
\]
The displayed parameters are identified by translation, so this is
the actual contact comparison.  It has order exactly two provided
\[
                  \alpha\kappa_0l_1+\beta\kappa_1l_0\ne0.
\]
This excludes one genuine affine linear equation in $r,s$:
its $r$ and $s$ coefficients are nonzero.
At the third point the derivative in the second-factor direction is
\[
                               rb_{20}+sb_{21}.
\]
Its nonvanishing excludes one more genuine linear equation and
ensures both smoothness at the third point and transversality to
the common tangent.  Since $k$ is infinite, all these finitely
many proper affine lines in the $(r,s)$-plane can be avoided.
This proves the existence of the required local type with five
independent linear conditions for every fixed $g,\eta$.

It remains to combine this local type with global smoothness and
the two nonvanishing source Hasse factors.  Let
$B=X\times\mathbf P^1$ parametrize $(R,\tau)$, where $\tau$ is
a common translated tangent direction.  The two value conditions
at $R,R+g$, their two directional-derivative conditions, and the
value condition at $R+\eta$ are algebraic linear conditions on
$H^0(X,L)$.  Their rank-five locus $\Omega\subset B$ is nonempty
by the construction above.  It is open and irreducible, since
$B$ is irreducible.  Their vector kernels form a rank-four bundle
on $\Omega$, whose projectivization $\mathcal T_{g,\eta}$ is therefore
an irreducible $\mathbf P^3$-bundle.

The point $(R,\tau)=((O,O),\text{second-factor tangent})$ also
belongs to $\Omega$: the first four conditions define the contact
$\mathbf P^4$ considered earlier, and the third-value functional
$l_\eta$ is nonzero there by
Corollary~\ref{cor:DE-A4-third-value-smooth}.
That corollary gives a point of this fibre whose source is globally
smooth and satisfies $\det S\,Q_H(S)\ne0$.
These conditions define a nonempty open of $\mathcal T_{g,\eta}$.
The local type constructed above likewise defines a nonempty open:
it is specified by nonvanishing first derivatives, a nonzero
contact-order-two coefficient, and a nonzero third derivative
in the marked tangent direction.
The two opens meet because $\mathcal T_{g,\eta}$ is irreducible.
Their intersection proves the lemma in full.
\end{proof}

\begin{remark}[Scope of the local-type construction]
\label{rem:DE-A4-uniform-triple-local-scope}
The lemma supplies a globally smooth sharp source with the required
three marked branches on the main rank-five incidence.
Exclusion of other identifications under the full quotient by $K$
is a further geometric condition.  It can be combined with this
lemma once it gives a nonempty open on the same incidence.
The proof makes no assertion that a finite list of small-kernel
witnesses controls arbitrary kernels.

The retained universal program
\path{agents/obstructions_crystals/code/uniform_triple_local_type_checks.py}
checks the rank-five minor, both tangent equations, the independent
first gradients, both nonzero coefficients of the simple-contact
condition, the third vertical derivative, and the standard cubic
tangent-line coefficient $x^4+x$.
Its author run passed in $0.0512757739998051$ seconds; its complete
output is retained alongside the source.
Independent complete program reads and author-output-preserving reruns
passed in $0.04299991099833278$ seconds (root) and
$0.05460318699624622$ seconds (monodromy).
The finite-exception and irreducible-incidence arguments are proved
above and are not inferred from this algebraic check.
\end{remark}
